% !PRINTED 365
\weberpage{365}
%<<CONT
then it follows that
\[
T_n S_{n-1} - S_n T_{n-1} = (-1)^n (T_0 S_1 - T_1 S_0),
\]
and so in particular

\weq{11}{P_n Q_{n-1} - Q_n P_{n-1} = (-1)^n.}

This formula, of which we shall yet make manifold application, shows that the numbers $P_n$, $Q_n$ have no common divisor, and hence that the convergents $P_n : Q_n$ cannot be reduced by cancellation.

If the continued fraction is a finite one, that is, if the closing number $x_{n-1} = a_{n-1}$ is an integer, then we carry the formation of the convergents no further than to
\[
\frac{P_n}{Q_n} = \frac{P_{n-1} a_{n-1} + P_{n-2}}{Q_{n-1} a_{n-1} + Q_{n-2}},
\]
and formula (2) shows that then the \gesperrt{last convergent agrees with the value of the continued fraction.}

\begin{center}
§. 118.

Solution of indeterminate equations with two unknowns.
\end{center}

The formulas last found lead to the solution of a problem that occurs in many applications:

\gesperrt{Let $\alpha$, $\beta$ be two given integers with no common divisor; two other integers are to be found which satisfy the condition:}

\weq{1}{\alpha y - \beta x = 1.}

This problem, if it is solvable at all, has infinitely many solutions, all of which can be derived from any one of them.

For if $x_0$, $y_0$ is a solution, so that

\weq{2}{\alpha y_0 - \beta x_0 = 1,}
then by subtraction of (1) and (2) it follows that
\[
\alpha (y - y_0) = \beta (x - x_0),
\]
from which it is to be concluded, since $\alpha$ has with $\beta$ no common divisor
%CONT>>
